\subsection{Frustration-free criteria require their hypotheses}
Gosset and Mozgunov give local gap thresholds and, in their gapless
one-dimensional periodic setting, an $O(n^{-2})$ open-chain upper bound
\cite{gossetMozgunov}. Lemm and Mozgunov give finite-size criteria with
boundaries \cite{lemmMozgunov}; their two-dimensional criterion uses a
minimum over rhomboidal patches, and general boundaries require the
corresponding intersected patch shapes. A square-box exponent alone does
not provide those patch or edge estimates. Masaoka, Soejima and Watanabe
obtain $z\geq2$ under their frustration-free correlation hypotheses,
including the extended correlations treated there
\cite{masaokaSoejimaWatanabe}. Lemm and Lucia prove inverse-square critical scaling, up to logarithmic
factors, without a ground-state correlation assumption
\cite[Theorem~1.4 and Remark~1.5]{lemmLucia}. Their result controls an
infimum over specified finite-domain families for one fixed positive
finite-range frustration-free interaction on graphs with the required
Euclidean embedding bounds. A square-box estimate alone does not give
that domain-family infimum, and periodic boundary conditions require the
additional boundary information stated there. These results do not imply
an unconditional exclusion interval for the square-box exponents of the
present setting.

The specified uniform local decomposition here is frustrated: an onsite
flag reward competes with interior bond penalties. For a particular box,
adding $4LI$ lets one regroup the flag energy into positive interior and
boundary terms. That boundary-dependent regrouping does not identify the
original interaction with a uniformly frustration-free family. No general
necessity of frustration for every small exponent is proved here.

